%%%%%%%%%%%%%%%%%%%%%%%%%%%%%%%%%%%%%%%%%
% The Legrand Orange Book
% LaTeX Template
% Version 2.4 (26/09/2018)
%
% This template was downloaded from:
% http://www.LaTeXTemplates.com
%
% Original author:
% Mathias Legrand (legrand.mathias@gmail.com) with modifications by:
% Vel (vel@latextemplates.com)
%
% License:
% CC BY-NC-SA 3.0 (http://creativecommons.org/licenses/by-nc-sa/3.0/)
%
% Compiling this template:
% This template uses biber for its bibliography and makeindex for its index.
% When you first open the template, compile it from the command line with the 
% commands below to make sure your LaTeX distribution is configured correctly:
%
% 1) pdflatex main
% 2) makeindex main.idx -s StyleInd.ist
% 3) biber main
% 4) pdflatex main x 2
%
% After this, when you wish to update the bibliography/index use the appropriate
% command above and make sure to compile with pdflatex several times 
% afterwards to propagate your changes to the document.
%
% This template also uses a number of packages which may need to be
% updated to the newest versions for the template to compile. It is strongly
% recommended you update your LaTeX distribution if you have any
% compilation errors.
%
% Important note:
% Chapter heading images should have a 2:1 width:height ratio,
% e.g. 920px width and 460px height.
%
%%%%%%%%%%%%%%%%%%%%%%%%%%%%%%%%%%%%%%%%%

%----------------------------------------------------------------------------------------
%	PACKAGES AND OTHER DOCUMENT CONFIGURATIONS
%----------------------------------------------------------------------------------------

\documentclass[11pt,fleqn]{book} % Default font size and left-justified equations

%%%%%%%%%%%%%%%%%%%%%%%%%%%%%%%%%%%%%%%%%
% The Legrand Orange Book
% Structural Definitions File
% Version 2.1 (26/09/2018)
%
% Original author:
% Mathias Legrand (legrand.mathias@gmail.com) with modifications by:
% Vel (vel@latextemplates.com)
% 
% This file was downloaded from:
% http://www.LaTeXTemplates.com
%
% License:
% CC BY-NC-SA 3.0 (http://creativecommons.org/licenses/by-nc-sa/3.0/)
%
%%%%%%%%%%%%%%%%%%%%%%%%%%%%%%%%%%%%%%%%%

%----------------------------------------------------------------------------------------
%	VARIOUS REQUIRED PACKAGES AND CONFIGURATIONS
%----------------------------------------------------------------------------------------

\usepackage{graphicx} % Required for including pictures
\graphicspath{{Pictures/}} % Specifies the directory where pictures are stored

\usepackage{lipsum} % Inserts dummy text

\usepackage{tikz} % Required for drawing custom shapes

\usepackage[english]{babel} % English language/hyphenation

\usepackage{enumitem} % Customize lists
\setlist{nolistsep} % Reduce spacing between bullet points and numbered lists

\usepackage{booktabs} % Required for nicer horizontal rules in tables

\usepackage{xcolor} % Required for specifying colors by name
%\definecolor{gray}{RGB}{243,102,25} % Define the orange color used for highlighting throughout the book

%----------------------------------------------------------------------------------------
%	MARGINS
%----------------------------------------------------------------------------------------

\usepackage{geometry} % Required for adjusting page dimensions and margins

\geometry{
	paper=a4paper, % Paper size, change to letterpaper for US letter size
	top=3cm, % Top margin
	bottom=3cm, % Bottom margin
	left=3cm, % Left margin
	right=3cm, % Right margin
	headheight=14pt, % Header height
	footskip=1.4cm, % Space from the bottom margin to the baseline of the footer
	headsep=10pt, % Space from the top margin to the baseline of the header
	%showframe, % Uncomment to show how the type block is set on the page
}

%----------------------------------------------------------------------------------------
%	FONTS
%----------------------------------------------------------------------------------------

\usepackage{avant} % Use the Avantgarde font for headings
%\usepackage{times} % Use the Times font for headings
\usepackage{mathptmx} % Use the Adobe Times Roman as the default text font together with math symbols from the Sym­bol, Chancery and Com­puter Modern fonts

\usepackage{microtype} % Slightly tweak font spacing for aesthetics
\usepackage[utf8]{inputenc} % Required for including letters with accents
\usepackage[T1]{fontenc} % Use 8-bit encoding that has 256 glyphs

%----------------------------------------------------------------------------------------
%	BIBLIOGRAPHY AND INDEX
%----------------------------------------------------------------------------------------

\usepackage[style=numeric,citestyle=numeric,sorting=nyt,sortcites=true,autopunct=true,babel=hyphen,hyperref=true,abbreviate=false,backref=true,backend=biber]{biblatex}
\addbibresource{bibliography.bib} % BibTeX bibliography file
\defbibheading{bibempty}{}

\usepackage{calc} % For simpler calculation - used for spacing the index letter headings correctly
\usepackage{makeidx} % Required to make an index
\makeindex % Tells LaTeX to create the files required for indexing

%----------------------------------------------------------------------------------------
%	MAIN TABLE OF CONTENTS
%----------------------------------------------------------------------------------------

\usepackage{titletoc} % Required for manipulating the table of contents

\contentsmargin{0cm} % Removes the default margin

% Part text styling (this is mostly taken care of in the PART HEADINGS section of this file)
\titlecontents{part}
	[0cm] % Left indentation
	{\addvspace{20pt}\bfseries} % Spacing and font options for parts
	{}
	{}
	{}

% Chapter text styling
\titlecontents{chapter}
	[1.25cm] % Left indentation
	{\addvspace{12pt}\large\sffamily\bfseries} % Spacing and font options for chapters
	{\color{gray!60}\contentslabel[\Large\thecontentslabel]{1.25cm}\color{gray}} % Formatting of numbered sections of this type
	{\color{gray}} % Formatting of numberless sections of this type
	{\color{gray!60}\normalsize\;\titlerule*[.5pc]{.}\;\thecontentspage} % Formatting of the filler to the right of the heading and the page number

% Section text styling
\titlecontents{section}
	[1.25cm] % Left indentation
	{\addvspace{3pt}\sffamily\bfseries} % Spacing and font options for sections
	{\contentslabel[\thecontentslabel]{1.25cm}} % Formatting of numbered sections of this type
	{} % Formatting of numberless sections of this type
	{\hfill\color{black}\thecontentspage} % Formatting of the filler to the right of the heading and the page number

% Subsection text styling
\titlecontents{subsection}
	[1.25cm] % Left indentation
	{\addvspace{1pt}\sffamily\small} % Spacing and font options for subsections
	{\contentslabel[\thecontentslabel]{1.25cm}} % Formatting of numbered sections of this type
	{} % Formatting of numberless sections of this type
	{\ \titlerule*[.5pc]{.}\;\thecontentspage} % Formatting of the filler to the right of the heading and the page number

% Figure text styling
\titlecontents{figure}
	[1.25cm] % Left indentation
	{\addvspace{1pt}\sffamily\small} % Spacing and font options for figures
	{\thecontentslabel\hspace*{1em}} % Formatting of numbered sections of this type
	{} % Formatting of numberless sections of this type
	{\ \titlerule*[.5pc]{.}\;\thecontentspage} % Formatting of the filler to the right of the heading and the page number

% Table text styling
\titlecontents{table}
	[1.25cm] % Left indentation
	{\addvspace{1pt}\sffamily\small} % Spacing and font options for tables
	{\thecontentslabel\hspace*{1em}} % Formatting of numbered sections of this type
	{} % Formatting of numberless sections of this type
	{\ \titlerule*[.5pc]{.}\;\thecontentspage} % Formatting of the filler to the right of the heading and the page number

%----------------------------------------------------------------------------------------
%	MINI TABLE OF CONTENTS IN PART HEADS
%----------------------------------------------------------------------------------------

% Chapter text styling
\titlecontents{lchapter}
	[0em] % Left indentation
	{\addvspace{15pt}\large\sffamily\bfseries} % Spacing and font options for chapters
	{\color{gray}\contentslabel[\Large\thecontentslabel]{1.25cm}\color{gray}} % Chapter number
	{}  
	{\color{gray}\normalsize\sffamily\bfseries\;\titlerule*[.5pc]{.}\;\thecontentspage} % Page number

% Section text styling
\titlecontents{lsection}
	[0em] % Left indentation
	{\sffamily\small} % Spacing and font options for sections
	{\contentslabel[\thecontentslabel]{1.25cm}} % Section number
	{}
	{}

% Subsection text styling (note these aren't shown by default, display them by searchings this file for tocdepth and reading the commented text)
\titlecontents{lsubsection}
	[.5em] % Left indentation
	{\sffamily\footnotesize} % Spacing and font options for subsections
	{\contentslabel[\thecontentslabel]{1.25cm}}
	{}
	{}

%----------------------------------------------------------------------------------------
%	HEADERS AND FOOTERS
%----------------------------------------------------------------------------------------

\usepackage{fancyhdr} % Required for header and footer configuration

\pagestyle{fancy} % Enable the custom headers and footers

\renewcommand{\chaptermark}[1]{\markboth{\sffamily\normalsize\bfseries\chaptername\ \thechapter.\ #1}{}} % Styling for the current chapter in the header
\renewcommand{\sectionmark}[1]{\markright{\sffamily\normalsize\thesection\hspace{5pt}#1}{}} % Styling for the current section in the header

\fancyhf{} % Clear default headers and footers
\fancyhead[LE,RO]{\sffamily\normalsize\thepage} % Styling for the page number in the header
\fancyhead[LO]{\rightmark} % Print the nearest section name on the left side of odd pages
\fancyhead[RE]{\leftmark} % Print the current chapter name on the right side of even pages
%\fancyfoot[C]{\thepage} % Uncomment to include a footer

\renewcommand{\headrulewidth}{0.5pt} % Thickness of the rule under the header

\fancypagestyle{plain}{% Style for when a plain pagestyle is specified
	\fancyhead{}\renewcommand{\headrulewidth}{0pt}%
}

% Removes the header from odd empty pages at the end of chapters
\makeatletter
\renewcommand{\cleardoublepage}{
\clearpage\ifodd\c@page\else
\hbox{}
\vspace*{\fill}
\thispagestyle{empty}
\newpage
\fi}

%----------------------------------------------------------------------------------------
%	THEOREM STYLES
%----------------------------------------------------------------------------------------

\usepackage{amsmath,amsfonts,amssymb,amsthm} % For math equations, theorems, symbols, etc

\newcommand{\intoo}[2]{\mathopen{]}#1\,;#2\mathclose{[}}
\newcommand{\ud}{\mathop{\mathrm{{}d}}\mathopen{}}
\newcommand{\intff}[2]{\mathopen{[}#1\,;#2\mathclose{]}}
\renewcommand{\qedsymbol}{$\blacksquare$}
\newtheorem{notation}{Notation}[chapter]

% Boxed/framed environments
\newtheoremstyle{graynumbox}% Theorem style name
{0pt}% Space above
{0pt}% Space below
{\normalfont}% Body font
{}% Indent amount
{\small\bf\sffamily\color{gray}}% Theorem head font
{\;}% Punctuation after theorem head
{0.25em}% Space after theorem head
{\small\sffamily\color{gray}\thmname{#1}\nobreakspace\thmnumber{\@ifnotempty{#1}{}\@upn{#2}}% Theorem text (e.g. Theorem 2.1)
\thmnote{\nobreakspace\the\thm@notefont\sffamily\bfseries\color{black}---\nobreakspace#3.}} % Optional theorem note

\newtheoremstyle{blacknumex}% Theorem style name
{5pt}% Space above
{5pt}% Space below
{\normalfont}% Body font
{} % Indent amount
{\small\bf\sffamily}% Theorem head font
{\;}% Punctuation after theorem head
{0.25em}% Space after theorem head
{\small\sffamily{\tiny\ensuremath{\blacksquare}}\nobreakspace\thmname{#1}\nobreakspace\thmnumber{\@ifnotempty{#1}{}\@upn{#2}}% Theorem text (e.g. Theorem 2.1)
\thmnote{\nobreakspace\the\thm@notefont\sffamily\bfseries---\nobreakspace#3.}}% Optional theorem note

\newtheoremstyle{blacknumbox} % Theorem style name
{0pt}% Space above
{0pt}% Space below
{\normalfont}% Body font
{}% Indent amount
{\small\bf\sffamily}% Theorem head font
{\;}% Punctuation after theorem head
{0.25em}% Space after theorem head
{\small\sffamily\thmname{#1}\nobreakspace\thmnumber{\@ifnotempty{#1}{}\@upn{#2}}% Theorem text (e.g. Theorem 2.1)
\thmnote{\nobreakspace\the\thm@notefont\sffamily\bfseries---\nobreakspace#3.}}% Optional theorem note

% Non-boxed/non-framed environments
\newtheoremstyle{graynum}% Theorem style name
{5pt}% Space above
{5pt}% Space below
{\normalfont}% Body font
{}% Indent amount
{\small\bf\sffamily\color{gray}}% Theorem head font
{\;}% Punctuation after theorem head
{0.25em}% Space after theorem head
{\small\sffamily\color{gray}\thmname{#1}\nobreakspace\thmnumber{\@ifnotempty{#1}{}\@upn{#2}}% Theorem text (e.g. Theorem 2.1)
\thmnote{\nobreakspace\the\thm@notefont\sffamily\bfseries\color{black}---\nobreakspace#3.}} % Optional theorem note
\makeatother

% Defines the theorem text style for each type of theorem to one of the three styles above
\newcounter{dummy} 
\numberwithin{dummy}{section}
\theoremstyle{graynumbox}
\newtheorem{theoremeT}[dummy]{Theorem}
\newtheorem{problem}{Problem}[chapter]
\newtheorem{exerciseT}{Exercise}[chapter]
\theoremstyle{blacknumex}
\newtheorem{exampleT}{Example}[chapter]
\theoremstyle{blacknumbox}
\newtheorem{vocabulary}{Vocabulary}[chapter]
\newtheorem{definitionT}{Definition}[section]
\newtheorem{corollaryT}[dummy]{Corollary}
\theoremstyle{graynum}
\newtheorem{proposition}[dummy]{Proposition}

%----------------------------------------------------------------------------------------
%	DEFINITION OF COLORED BOXES
%----------------------------------------------------------------------------------------

\RequirePackage[framemethod=default]{mdframed} % Required for creating the theorem, definition, exercise and corollary boxes

% Theorem box
\newmdenv[skipabove=7pt,
skipbelow=7pt,
backgroundcolor=black!5,
linecolor=gray,
innerleftmargin=5pt,
innerrightmargin=5pt,
innertopmargin=5pt,
leftmargin=0cm,
rightmargin=0cm,
innerbottommargin=5pt]{tBox}

% Exercise box	  
\newmdenv[skipabove=7pt,
skipbelow=7pt,
rightline=false,
leftline=true,
topline=false,
bottomline=false,
backgroundcolor=gray!10,
linecolor=gray,
innerleftmargin=5pt,
innerrightmargin=5pt,
innertopmargin=5pt,
innerbottommargin=5pt,
leftmargin=0cm,
rightmargin=0cm,
linewidth=4pt]{eBox}	

% Definition box
\newmdenv[skipabove=7pt,
skipbelow=7pt,
rightline=false,
leftline=true,
topline=false,
bottomline=false,
linecolor=gray,
innerleftmargin=5pt,
innerrightmargin=5pt,
innertopmargin=0pt,
leftmargin=0cm,
rightmargin=0cm,
linewidth=4pt,
innerbottommargin=0pt]{dBox}	

% Corollary box
\newmdenv[skipabove=7pt,
skipbelow=7pt,
rightline=false,
leftline=true,
topline=false,
bottomline=false,
linecolor=gray,
backgroundcolor=black!5,
innerleftmargin=5pt,
innerrightmargin=5pt,
innertopmargin=5pt,
leftmargin=0cm,
rightmargin=0cm,
linewidth=4pt,
innerbottommargin=5pt]{cBox}

% Creates an environment for each type of theorem and assigns it a theorem text style from the "Theorem Styles" section above and a colored box from above
\newenvironment{theorem}{\begin{tBox}\begin{theoremeT}}{\end{theoremeT}\end{tBox}}
\newenvironment{exercise}{\begin{eBox}\begin{exerciseT}}{\hfill{\color{gray}\tiny\ensuremath{\blacksquare}}\end{exerciseT}\end{eBox}}				  
\newenvironment{definition}{\begin{dBox}\begin{definitionT}}{\end{definitionT}\end{dBox}}	
\newenvironment{example}{\begin{exampleT}}{\hfill{\tiny\ensuremath{\blacksquare}}\end{exampleT}}		
\newenvironment{corollary}{\begin{cBox}\begin{corollaryT}}{\end{corollaryT}\end{cBox}}	

%----------------------------------------------------------------------------------------
%	REMARK ENVIRONMENT
%----------------------------------------------------------------------------------------

\newenvironment{remark}{\par\vspace{10pt}\small % Vertical white space above the remark and smaller font size
\begin{list}{}{
\leftmargin=35pt % Indentation on the left
\rightmargin=25pt}\item\ignorespaces % Indentation on the right
\makebox[-2.5pt]{\begin{tikzpicture}[overlay]
\node[draw=gray!60,line width=1pt,circle,fill=gray!25,font=\sffamily\bfseries,inner sep=2pt,outer sep=0pt] at (-15pt,0pt){\textcolor{gray}{R}};\end{tikzpicture}} % Orange R in a circle
\advance\baselineskip -1pt}{\end{list}\vskip5pt} % Tighter line spacing and white space after remark

%----------------------------------------------------------------------------------------
%	SECTION NUMBERING IN THE MARGIN
%----------------------------------------------------------------------------------------

\makeatletter
\renewcommand{\@seccntformat}[1]{\llap{\textcolor{gray}{\csname the#1\endcsname}\hspace{1em}}}                    
\renewcommand{\section}{\@startsection{section}{1}{\z@}
{-4ex \@plus -1ex \@minus -.4ex}
{1ex \@plus.2ex }
{\normalfont\large\sffamily\bfseries}}
\renewcommand{\subsection}{\@startsection {subsection}{2}{\z@}
{-3ex \@plus -0.1ex \@minus -.4ex}
{0.5ex \@plus.2ex }
{\normalfont\sffamily\bfseries}}
\renewcommand{\subsubsection}{\@startsection {subsubsection}{3}{\z@}
{-2ex \@plus -0.1ex \@minus -.2ex}
{.2ex \@plus.2ex }
{\normalfont\small\sffamily\bfseries}}                        
\renewcommand\paragraph{\@startsection{paragraph}{4}{\z@}
{-2ex \@plus-.2ex \@minus .2ex}
{.1ex}
{\normalfont\small\sffamily\bfseries}}

%----------------------------------------------------------------------------------------
%	PART HEADINGS
%----------------------------------------------------------------------------------------

% Numbered part in the table of contents
\newcommand{\@mypartnumtocformat}[2]{%
	\setlength\fboxsep{0pt}%
	\noindent\colorbox{gray!20}{\strut\parbox[c][.7cm]{\ecart}{\color{gray!70}\Large\sffamily\bfseries\centering#1}}\hskip\esp\colorbox{gray!40}{\strut\parbox[c][.7cm]{\linewidth-\ecart-\esp}{\Large\sffamily\centering#2}}%
}

% Unnumbered part in the table of contents
\newcommand{\@myparttocformat}[1]{%
	\setlength\fboxsep{0pt}%
	\noindent\colorbox{gray!40}{\strut\parbox[c][.7cm]{\linewidth}{\Large\sffamily\centering#1}}%
}

\newlength\esp
\setlength\esp{4pt}
\newlength\ecart
\setlength\ecart{1.2cm-\esp}
\newcommand{\thepartimage}{}%
\newcommand{\partimage}[1]{\renewcommand{\thepartimage}{#1}}%
\def\@part[#1]#2{%
\ifnum \c@secnumdepth >-2\relax%
\refstepcounter{part}%
\addcontentsline{toc}{part}{\texorpdfstring{\protect\@mypartnumtocformat{\thepart}{#1}}{\partname~\thepart\ ---\ #1}}
\else%
\addcontentsline{toc}{part}{\texorpdfstring{\protect\@myparttocformat{#1}}{#1}}%
\fi%
\startcontents%
\markboth{}{}%
{\thispagestyle{empty}%
\begin{tikzpicture}[remember picture,overlay]%
\node at (current page.north west){\begin{tikzpicture}[remember picture,overlay]%	
\fill[gray!20](0cm,0cm) rectangle (\paperwidth,-\paperheight);
\node[anchor=north] at (4cm,-3.25cm){\color{gray!40}\fontsize{220}{100}\sffamily\bfseries\thepart}; 
\node[anchor=south east] at (\paperwidth-1cm,-\paperheight+1cm){\parbox[t][][t]{8.5cm}{
\printcontents{l}{0}{\setcounter{tocdepth}{1}}% The depth to which the Part mini table of contents displays headings; 0 for chapters only, 1 for chapters and sections and 2 for chapters, sections and subsections
}};
\node[anchor=north east] at (\paperwidth-1.5cm,-3.25cm){\parbox[t][][t]{15cm}{\strut\raggedleft\color{white}\fontsize{30}{30}\sffamily\bfseries#2}};
\end{tikzpicture}};
\end{tikzpicture}}%
\@endpart}
\def\@spart#1{%
\startcontents%
\phantomsection
{\thispagestyle{empty}%
\begin{tikzpicture}[remember picture,overlay]%
\node at (current page.north west){\begin{tikzpicture}[remember picture,overlay]%	
\fill[gray!20](0cm,0cm) rectangle (\paperwidth,-\paperheight);
\node[anchor=north east] at (\paperwidth-1.5cm,-3.25cm){\parbox[t][][t]{15cm}{\strut\raggedleft\color{white}\fontsize{30}{30}\sffamily\bfseries#1}};
\end{tikzpicture}};
\end{tikzpicture}}
\addcontentsline{toc}{part}{\texorpdfstring{%
\setlength\fboxsep{0pt}%
\noindent\protect\colorbox{gray!40}{\strut\protect\parbox[c][.7cm]{\linewidth}{\Large\sffamily\protect\centering #1\quad\mbox{}}}}{#1}}%
\@endpart}
\def\@endpart{\vfil\newpage
\if@twoside
\if@openright
\null
\thispagestyle{empty}%
\newpage
\fi
\fi
\if@tempswa
\twocolumn
\fi}

%----------------------------------------------------------------------------------------
%	CHAPTER HEADINGS
%----------------------------------------------------------------------------------------

% A switch to conditionally include a picture, implemented by Christian Hupfer
\newif\ifusechapterimage
\usechapterimagetrue
\newcommand{\thechapterimage}{}%
\newcommand{\chapterimage}[1]{\ifusechapterimage\renewcommand{\thechapterimage}{#1}\fi}%
\newcommand{\autodot}{.}
\def\@makechapterhead#1{%
{\parindent \z@ \raggedright \normalfont
\ifnum \c@secnumdepth >\m@ne
\if@mainmatter
\begin{tikzpicture}[remember picture,overlay]
\node at (current page.north west)
{\begin{tikzpicture}[remember picture,overlay]
\node[anchor=north west,inner sep=0pt] at (0,0) {\ifusechapterimage\includegraphics[width=\paperwidth]{\thechapterimage}\fi};
\draw[anchor=west] (\Gm@lmargin,-9cm) node [line width=2pt,rounded corners=15pt,draw=gray,fill=white,fill opacity=0.5,inner sep=15pt]{\strut\makebox[22cm]{}};
\draw[anchor=west] (\Gm@lmargin+.3cm,-9cm) node {\huge\sffamily\bfseries\color{black}\thechapter\autodot~#1\strut};
\end{tikzpicture}};
\end{tikzpicture}
\else
\begin{tikzpicture}[remember picture,overlay]
\node at (current page.north west)
{\begin{tikzpicture}[remember picture,overlay]
\node[anchor=north west,inner sep=0pt] at (0,0) {\ifusechapterimage\includegraphics[width=\paperwidth]{\thechapterimage}\fi};
\draw[anchor=west] (\Gm@lmargin,-9cm) node [line width=2pt,rounded corners=15pt,draw=gray,fill=white,fill opacity=0.5,inner sep=15pt]{\strut\makebox[22cm]{}};
\draw[anchor=west] (\Gm@lmargin+.3cm,-9cm) node {\huge\sffamily\bfseries\color{black}#1\strut};
\end{tikzpicture}};
\end{tikzpicture}
\fi\fi\par\vspace*{270\p@}}}

%-------------------------------------------

\def\@makeschapterhead#1{%
\begin{tikzpicture}[remember picture,overlay]
\node at (current page.north west)
{\begin{tikzpicture}[remember picture,overlay]
\node[anchor=north west,inner sep=0pt] at (0,0) {\ifusechapterimage\includegraphics[width=\paperwidth]{\thechapterimage}\fi};
\draw[anchor=west] (\Gm@lmargin,-9cm) node [line width=2pt,rounded corners=15pt,draw=gray,fill=white,fill opacity=0.5,inner sep=15pt]{\strut\makebox[22cm]{}};
\draw[anchor=west] (\Gm@lmargin+.3cm,-9cm) node {\huge\sffamily\bfseries\color{black}#1\strut};
\end{tikzpicture}};
\end{tikzpicture}
\par\vspace*{270\p@}}
\makeatother

%----------------------------------------------------------------------------------------
%	LINKS
%----------------------------------------------------------------------------------------

\usepackage{hyperref}
\hypersetup{hidelinks,backref=true,pagebackref=true,hyperindex=true,colorlinks=false,breaklinks=true,urlcolor=gray,bookmarks=true,bookmarksopen=false}

\usepackage{bookmark}
\bookmarksetup{
open,
numbered,
addtohook={%
\ifnum\bookmarkget{level}=0 % chapter
\bookmarksetup{bold}%
\fi
\ifnum\bookmarkget{level}=-1 % part
\bookmarksetup{color=gray,bold}%
\fi
}
}
 % Insert the commands.tex file which contains the majority of the structure behind the template

%\hypersetup{pdftitle={Title},pdfauthor={Author}} % Uncomment and fill out to include PDF metadata for the author and title of the book

\usepackage[]{nicefrac}
\usepackage[]{wrapfig}
\usepackage[]{bm}

\usetikzlibrary{arrows}
\usetikzlibrary{shapes.geometric, math, positioning, calc, patterns, angles, quotes}
\usetikzlibrary{patterns.meta,decorations.pathmorphing}

\newcommand{\semaphore}[3]{% #1: color of circle,
                           % #2: color of semicircle
                           % #3: angle of semicircle 
  \tikz[node distance=0mm,baseline]
       {
         \node (s1) [circle, fill=#1, minimum size=6mm] {};
         \node      [semicircle, fill=#2, 
           inner sep=0pt, outer sep=0pt, minimum size=3mm,
           anchor=south,
           at={(s1.center)}, rotate=#3] {};
       }
}

\usepackage[]{circuitikz}
\usetikzlibrary{lindenmayersystems}


\usepackage[]{subfigure}

%----------------------------------------------------------------------------------------

\begin{document}


%----------------------------------------------------------------------------------------
%	TITLE PAGE
%----------------------------------------------------------------------------------------

\begingroup
\thispagestyle{empty} % Suppress headers and footers on the title page
\begin{tikzpicture}[remember picture,overlay]
\node[inner sep=0pt] (background) at (current page.center) {\includegraphics[width=\paperwidth]{background.pdf}};
\draw (current page.center) node [fill=gray!30!white,fill opacity=0.6,text opacity=1,inner sep=1cm]{\Huge\centering\bfseries\sffamily\parbox[c][][t]{\paperwidth}{\centering Nonlinear dynamics\\[15pt] % Book title
{\Large A Profound Subtitle}\\[20pt] % Subtitle
{\huge Dr. Jean-Christophe Loiseau}}}; % Author name
\end{tikzpicture}
\vfill
\endgroup

%----------------------------------------------------------------------------------------
%	COPYRIGHT PAGE
%----------------------------------------------------------------------------------------

\newpage
~\vfill
\thispagestyle{empty}

\noindent Copyright \copyright\ 2019 John Smith\\ % Copyright notice

\noindent \textsc{Published by Publisher}\\ % Publisher

\noindent \textsc{book-website.com}\\ % URL

\noindent Licensed under the Creative Commons Attribution-NonCommercial 3.0 Unported License (the ``License''). You may not use this file except in compliance with the License. You may obtain a copy of the License at \url{http://creativecommons.org/licenses/by-nc/3.0}. Unless required by applicable law or agreed to in writing, software distributed under the License is distributed on an \textsc{``as is'' basis, without warranties or conditions of any kind}, either express or implied. See the License for the specific language governing permissions and limitations under the License.\\ % License information, replace this with your own license (if any)

\noindent \textit{First printing, March 2019} % Printing/edition date

%----------------------------------------------------------------------------------------
%	TABLE OF CONTENTS
%----------------------------------------------------------------------------------------

%\usechapterimagefalse % If you don't want to include a chapter image, use this to toggle images off - it can be enabled later with \usechapterimagetrue

\chapterimage{chapter_head_1.pdf} % Table of contents heading image

\pagestyle{empty} % Disable headers and footers for the following pages

\tableofcontents % Print the table of contents itself

\cleardoublepage % Forces the first chapter to start on an odd page so it's on the right side of the book

\pagestyle{fancy} % Enable headers and footers again

\chapter*{Introduction}

\lipsum[7]


%----------------------------------------------------------------------------------------
%	PART ONE : DYNAMICAL SYSTEMS
%----------------------------------------------------------------------------------------

\part{Introduction to dynamical systems}

% --> First-order systems
\chapter{First-order systems}\label{chap: first-order systems}

\section{Dynamical systems on the real number line}\label{sec: dyn sys on the real number line}

\subsection{List of problems}

\begin{problem}
  Consider the simple resistor-capacitor circuit shown in figure~\ref{fig: RC circuit}.
  The equation governing the evolution of the voltage $Q(t)$ accross the capacitor is given by
  %
  \[
  \dot{Q} = - \dfrac{Q}{RC} + \dfrac{V_0}{RC}
  \]
  %
  where $R$ is the resistance, $C$ is the capacitance and $V_0 > 0$ the voltage of the power supply.
  For $t < 0$, the voltage accross the capacitor is assumed to be $0$.
  At $t = 0$, the switch is closed and current starts to flow in the circuit.

  \bigskip

  \begin{enumerate}
  \item[a)] Show graphically that the system has a single fixed point.
  \item[b)] Show that the equation above can be recast as
    %
    \[
    \dfrac{dx}{d\tau} = 1 - x
    \]
    %
    if one introduces the rescaled variable $Q(t) = V_0 x(t)$ and time $t = RC \tau$.
  \item[c)] Solve the equation analytically and show that
    %
    \[
    \lim_{\tau \to +\infty} x(\tau) = 1
    \]
    %
    and give the expression for $Q(t)$.
  \end{enumerate}
\end{problem}

\begin{figure}
  \centering
  \begin{circuitikz}
    \draw (0, 0) to[battery1, invert] (0, 3) to [nos] (2.5, 3);
    \draw (2.5, 3) to [R, i>^=$i(t)$, l=$R$] (2.5, 0);
    \draw (2.5, 0) to [C, l=$C$] (0, 0);
  \end{circuitikz}
  \caption{RC circuit for Problem 1.1}\label{fig: RC circuit}
\end{figure}

\bigskip

\begin{problem}
  Consider the following population growth model
  %
  \[
  \dot{P} = \mu P \left( 1 - \dfrac{P}{N} \right)
  \]
  %
  known as the logistic equation (also sometimes called the Verhulst model after Pierre François Verhulst who proposed it during the 1840's).
  It can be used to model the growth of a population in an environment with limited food supply such as fish in a small lake.
  In that context, $P(t)$ denotes the current population size, $\mu > 0$ its reproduction rate and $N$ is the \emph{carrying capacity}, i.e.\ the maximum population size that can be sustained in that specific environment.

  \bigskip

  \begin{enumerate}
  \item[a)] Introducing the rescaled population size $x(t)$ such that $P(t) = N x(t)$, show that the governing equation for $x(t)$ is
    %
    \[
    \dot{x} = \mu x \left( 1 - x \right).
    \]

  \item[b)] Sketch the phase line of the system assuming $\mu > 0$.
    How does the population evolves if $x(0)$ is only slightly larger than 0 ?
    Same question for $x(0)$ much larger than unity.

  \item[c)] This logistic model belongs to the class of \emph{Bernoulli equations} which can be solved analytically.
    Introducing the change of variable $y = \nicefrac{1}{x}$, show that the equation for $y(t)$ is linear.

  \item[d)] Solve the equation for $y(t)$ and thus for $x(t)$.
    Are the predictions of your solution for $x(0) \simeq 0$ and $x(0) \gg 1$ consistent with your intuition gained from the phase line ?
  \end{enumerate}
\end{problem}

\section{Elementary bifurcations}\label{sec: elementary bifurcations}

\subsection{List of problems}

\begin{problem}
  Consider the following system
  %
  \[
  \dot{x} = -x \left( x^2 - 2x - \mu \right)
  \]
  %
  where $x \in \mathbb{R}$ is the state of the system and $\mu \in \mathbb{R}$ is our control parameter.

  \bigskip

  \begin{enumerate}
  \item[a)] Compute all the branches of solutions and discuss their existence and stability properties.
  \item[b)] Show that a saddle-node bifurcation occurs at $\mu = -1$ and a transcritical one at $\mu = 0$.
  \item[c)] Sketch the bifurcation diagram of the system.
  \end{enumerate}
\end{problem}

\bigskip

\begin{problem}
  Consider a simple pendulum of mass $m$ and length $L$ in the gravitational field of earth ($g$) and driven by a constant torque $\Gamma$.
  Starting from Newton's principles, the equation of motion reads
  %
  \[
  m L^2 \ddot{\theta} + b \dot{\theta} + mgL \sin(\theta) = \Gamma.
  \]
  %
  Although it is a second-order equation (i.e.\ it involves a second derivative with respect to time), we'll see that under certain assumptions about the system it can be approximated by a first-order equation.
  Our aim will be to determine the type of bifurcation experienced by the system as the applied torque $\Gamma$ is varied.

  \bigskip

  \begin{enumerate}
  \item[a)] Introducing a time scale $\tau$ such that $t \mapsto \tau t$, show that we have the following choices for $\tau$ :
    %
    \[
    \tau = \sqrt{\dfrac{L}{g}} \quad \text{and} \quad \tau = \dfrac{b}{mgL}
    \]
    %
    and discuss their physical meaning.

  \item[b)] If we are in the over-damped situation, the equation of motion can be reduced to
    %
    \[
    \dot{\theta}  = \gamma - \sin(\theta)
    \]
    %
    where $\gamma = \nicefrac{\Gamma}{mgL}$.
    Sketch the phase line of the system for different values of the control parameter $\gamma$.

  \item[c)] As $\gamma \to 1$ from above, a meta-stable fixed point is created at $\theta = \nicefrac{\pi}{2}$.
    Introducing $\mu = \gamma - 1$ and using a second-order Taylor expansion of the equation in the vicinity of this fixed point, show that the dynamics of a small perturbation can be approximated by
    %
    \[
    \dot{\eta} = \mu + \dfrac{1}{2}\eta^2.
    \]
    %
    What type of bifurcation is this ?

  \item[d)] For $\vert \gamma \vert < 1$, the system admits two equilibrium solutions.
    Discuss their stability properties and their physical interpretation.
  \end{enumerate}
\end{problem}

\bigskip

\begin{problem}
  Consider the following mechanical system : an inverted pendulum made of a massless rigid rod is pivoted about its lower end with a torsion spring providing a restoring torque proportional to the angular displacement from the vertical equilibrium.
  A load, in the form of an attached mass, is applied vertically at the top of the rod.
  Starting from Newton's principle, the equation of motion is
  %
  \[
  I \ddot{\theta} + \beta \dot{\theta} + k \theta - PL \sin(\theta) = 0
  \]
  %
  where $I$ the moment of inertia of the system, $k$ is the torsional spring constant, $P$ is the applied load, $L$ is the length of the rod and $\beta$ is the damping coefficient.
  Our goal will be to determine when the system bifurcates and what type of bifurcation it experiences as the applied load $P$ is varied.

  \bigskip

  \begin{enumerate}

  \item[a)] Show that the equation of motion is invariant with respect to the transformation $\theta \mapsto -\theta$.
    What does this invariance tells you about the properties of the system ?

  \item[b)] As before, we'll consider the over-damped situation (i.e. friction dominates).
    Once time has been rescaled using $\tau = \nicefrac{\beta}{PL}$, the equation of motion reduces to
    %
    \[
    \dot{\theta} + \dfrac{k}{PL}\theta - \sin(\theta) = 0.
    \]
    %
    Using a third-order Taylor expansion of $\sin(\theta)$ around $\theta = 0$, show that the system experiences a bifurcation for $\nicefrac{k}{PL} = 1$.
    Sketch the bifurcation diagram.
    What type of bifurcation is this ?

  \item[c)] In a realistic system, one may have some imperfections breaking the symmetry.
    Redo the same analysis for the perturbed system
    %
    \[
    I \ddot{\theta} + \beta \dot{\theta} + k \theta - PL \left( h + \sin(\theta) \right) = 0
    \]
    %
    where $h$ models a small imbalance in the applied load.
    You can assume that $h > 0$.
    How does the bifurcation change ?
  \end{enumerate}

\end{problem}

\bigskip

\begin{problem}
  Let consider once more the evolution of a fish population model by the following logistic growth equation
  %
  \[
  \dot{P} = r P \left( 1 - \dfrac{P}{N} \right)
  \]
  %
  where $P(t) \in \mathbb{R}$ is the fish population, $r$ is the growth rate due to reproduction and $N$ is the carrying capacity of the environment.
  In this exercise, we'll investigate the influence of different harvesting strategies on the long time evolution of the fish population.
  Two different strategies will be considered.

  \paragraph{Constant harvesting}
  The first strategy considered is constant harvesting.
  In this strategy, fish are harvested at a constant rate irrespective of the current population level.
  If we denote by $H$ this rate, our model becomes
  %
  \[
  \dot{P} = r P \left( 1 - \dfrac{P}{N} \right) - H
  \]
  %
  supplemented with the initial condition $P(0) = P_0$.

  \bigskip

  \begin{enumerate}
  \item[a)] The model above has three parameters : the reproduction rate $r$, the carrying capacity $N$ and the harvesting rate $H$.
    Using a suitable rescaling, show that one of the parameter can be eliminated so that our model can be recast as
    %
    \[
    \dot{x} = r x \left( 1 - x \right) - h.
    \]

  \item[b)] Assuming $r$ is a positive constant, compute the fixed points of the model and study their linear stability properties.

  \item[c)] Sketch the bifurcation diagram.
    What type of bifurcation is this ?
    From a biological point of view, what happens to the fish population when $h$ is too large ?

  \end{enumerate}

  \paragraph{Proportional harvesting}
  In this second strategy, fish are harvested at a rate proportional to the current population level.
  Our population model then becomes
  %
  \[
  \dot{P} = r P \left( 1 - \dfrac{P}{N} \right) - HP.
  \]

  \begin{enumerate}
  \item[a)] As for the previous model, show that the current one can be written as
    %
    \[
    \dot{x} = r x \left( 1 - x \right) - hx
    \]
    %
    using a suitable rescaling.

  \item[b)] Assuming $r$ is a positive constant, compute the fixed points of the model and study their linear stability properties.

  \item[c)] Sketch the bifurcation diagram.
    What type of bifurcation is this ?
    From a biological point of view, what happens to the fish population when $h$ is too large ?
    How is this different from the previous strategy ?

  \end{enumerate}
\end{problem}


% --> Second-order systems
\chapter{Second-order systems}\label{chap: second-order systems}

\section{Linear systems}\label{sec: linear systems}

\subsection{List of problems}

\subsubsection*{Linear algebra}

\begin{problem}
  Given a matrix $\bm{A} \in \mathbb{R}^{2 \times 2}$, show that its characteristic polynomial can be expressed as
  %
  \[
  \lambda^2 - \text{tr}(\bm{A}) \lambda + \det(\bm{A}) = 0
  \]
  %
  where $\text{tr}(\bm{A})$ denotes the trace of the matrix and $\det(\bm{A})$ its determinant.
\end{problem}

\bigskip

\begin{problem}
  Show that, if $\bm{A} \in \mathbb{R}^{n \times n}$ is a symmetric matrix (i.e.\ $\bm{A}^T = \bm{A}$), then all of its eigenvalues $\lambda$ are real and its eigenvectors $\bm{v}$'s form an orthonormal basis for $\mathbb{R}^n$.
\end{problem}

\bigskip

\begin{problem}
  Show that if $\bm{A}$ is a symmetric negative definite matrix (i.e.\ $\bm{A}^T = \bm{A}$ and all of its eigenvalues are negative) then
  %
  \[
  \dot{\bm{x}} = \bm{Ax}
  \]
  %
  implies
  %
  \[
  \dfrac{1}{2} \dfrac{d}{dt} \| \bm{x} \|_2^2 = - \| \bm{Rx} \|_2^2
  \]
  %
  where $\bm{R}$ is the Cholesky factor of $-\bm{A}$ (i.e.\ $\bm{A} = - \bm{R}^T \bm{R}$ where $\bm{R}^T \bm{R}$ is a symmetric positive definite matrix).
  What does the equation above implies for the evolution of the energy in the system ?
\end{problem}

\bigskip

\begin{problem}
  Show that any matrix of the form
  %
  \[
  \bm{A} = \begin{bmatrix} \lambda & a \\ 0 & \lambda \end{bmatrix}
  \]
  %
  with $a \neq 0$ has only a one-dimensional eigenspace corresponding to the eigenvalue $\lambda$.
  Then solve the system $\dot{\bm{x}} = \bm{Ax}$ and sketch the phase portrait.
\end{problem}

\bigskip

\begin{problem}
  Show that the eigenvalues of a skew-symmetric matrix
  %
  \[
  \bm{A} = \begin{bmatrix} 0 & a \\ -a & 0 \end{bmatrix}
  \]
  %
  are complex conjugate imaginary numbers.
\end{problem}

\bigskip

\begin{problem}
  Consider the linear dynamical system $\dot{\bm{x}} = \bm{Ax}$ given by
  %
  \[
  \dfrac{d}{dt} \begin{bmatrix} x \\ y \end{bmatrix}
  =
  \begin{bmatrix}
    1 & -1 \\ 1 & 1
  \end{bmatrix}
  \begin{bmatrix}
    x \\ y
  \end{bmatrix}.
  \]

  \begin{enumerate}
  \item[a)] Show that its eigenvalues are given by $\lambda = 1 \pm i$ and compute its eigenvectors $\bm{v}_1$ and $\bm{v}_2$.

  \item[b)] The general solution is $\bm{x}(t) = \alpha e^{\lambda_1 t} \bm{v}_1 + \beta e^{\lambda_2t} \bm{v}_2$.
    If $\bm{x}(0) \in \mathbb{R}^2$ then $\bm{x}(t) \in \mathbb{R}^2 \ \forall t$.
    In such a case, express $\bm{x}(t)$ purely in terms of real-valued functions.

  \item[c)] Which aspect of the dynamics do the real and imaginary parts of $\lambda$ characterize ?
  \end{enumerate}
\end{problem}

\subsubsection*{Harmonic oscillator (with and without damping)}

\begin{problem}
  Consider the canonical harmonic oscillator whose equation of motion is given by
  %
  \[
  \ddot{x} + \omega^2 x = 0.
  \]

  \begin{enumerate}
  \item[a)] Introducing a suitable potential $V(x) : \mathbb{R} \to \mathbb{R}$, show that the equation for the harmonic oscillator can be rewritten as
    %
    \[
    \ddot{x} + \dfrac{dV}{dx} = 0.
    \]

  \item[b)] Show that the equation above implies
    %
    \[
    \dfrac{d}{dt} \left( \dfrac{1}{2} \dot{x}^2 + V(x) \right) = 0.
    \]
    %
    What is the meaning of this equation ?
    What does the quantity $H(x, \dot{x}) = \dfrac{1}{2} \dot{x}^2 + V(x)$ represent from a physical point of view ?

  \item[c)] Given the expression for the potential $V(x)$ found in the previous question, show that the trajectories of the harmonic oscillator in phase space correspond to ellipses of the form $\omega^2 x^2 + y^2 = C$ with $C > 0$.
  \end{enumerate}
\end{problem}

\bigskip

\begin{problem}
  Consider the simple mechanical system shown in figure~\ref{fig: mass-spring-damper}.
  Starting from Newton's principles, the equation of motion is given by
  %
  \[
  \ddot{x} + 2k \dot{x} + \omega_0^2 x = 0
  \]
  %
  where $x$ denotes the deviation from the equilibrium position, $k > 0$ is the friction coefficient and $\omega_0$ is the natural frequency of the system in the absence of friction.
  In this exercise, we aim at classifying all of the possible dynamics exhibited by the system as we vary its parameters.
  
  \bigskip
  
  \begin{enumerate}
  \item[a)] In its current form, our model depends explicitely on two parameters, namely the friction coefficient $k$ and the natural frequency of oscillation $\omega_0$.
    Using a suitable rescaling of time, show that the model can be rewritten as
    %
    \[
    \ddot{x} + 2 \mu \dot{x} + x = 0
    \]
    %
    and give the expression of $\mu$.
    
  \item[b)] Rewrite the model above as a two-dimensional linear system of the form $\dot{\bm{x}} = \bm{Ax}$.
    
  \item[c)] Compute the characteristic polynomial of the matrix $\bm{A}$ and classify the type of fixed point for all the possible cases.
    For all cases, sketch the phase portrait of the system.
    
  \item[d)] How do these results relate to the standard notions of overdamped, critically damped and underdamped oscillations ?
  \end{enumerate}
  
\end{problem}

\begin{figure}
  \centering
  \begin{tikzpicture}[>=stealth]
      \path[pattern={Lines[angle=45,distance={8pt/sqrt(2)}]}] (-2, 5) edge ++(4,0) rectangle ++ (4, 0.5);

      \draw[decorate,decoration={coil, segment length=5pt, aspect=0.7, amplitude=4pt,
          pre=lineto, pre length=5mm, post=lineto, post length=5mm}] (0.5, 5) -- (0.5, 2.5)
      node[left, draw, minimum size=1cm, fill=blue!20] (m) {$m$};

      \draw[] (-0.5, 5) -- (-0.5, 4) {};
      \draw[] (-0.25, 4) -- (-0.75, 4) {};

      \draw[] (-0.5, 3.75) -- (-0.5, 2.5) {};
      \draw[] (-0.25, 3.75) -- (-0.75, 3.75) {};

      \draw[] (-0.25, 3.75) -- (-0.25, 4.25) {};
      \draw[] (-0.75, 3.75) -- (-0.75, 4.25) {};

      \draw[blue] (m.center-|0, 0) ++ (0, 1) -- ++ (-2, 0)
      edge[<-,edge label'=$x$,shorten >=1pt] (m.center-|-2, 0)
      (m.center-|0, 0) ++ (0, -1) -- ++ (-2, 0) 
      edge[<-,edge label=$-x$,shorten >=1pt] (m.center-|-2, 0)
      (m.east) edge[dashed] (m.east-|2, 0) 
      (m.east-|2, 0) node[right] {Equilibrium};
  \end{tikzpicture}
  \caption{Mechanical system considered in Problem 2.2}\label{fig: mass-spring-damper}
\end{figure}

\bigskip

\begin{problem}
  Consider the electric circuit shown in figure~\ref{fig: RLC circuit}.
  The Kirchoff's voltage law implies
  %
  \[
  V_R(t) + V_L(t) + V_C(t) = V_{\text{in}}(t)
  \]
  %
  where $V_R$, $V_L$ and $V_C$ are the voltages accross R, L, and C respectively and $V_{\text{in}}$ is the (possibly time-varying) voltage from the source.
  Using Ohm's law, the voltage accross the resistor can be expressed as
  %
  \[
  V_R(t) = R i(t).
  \]
  %
  Similarly, the voltage $V_L$ accross the inductor is related to the intensity of the current by
  %
  \[
  V_L(t) = L \dfrac{di}{dt}.
  \]
  %
  Finally, the intensity $i(t)$ of the current is related to the voltage accross the capacitor by
  %
  \[
  i = C \dfrac{dV_C}{dt}.
  \]
  %
  Combining all of these consitutive equations together yields the following second-order differential equation
  %
  \[
  \dfrac{d^2 V_C}{dt^2} + \dfrac{R}{L} \dfrac{d V_c}{dt} + \dfrac{V_C}{LC} = \dfrac{V_{\text{in}}}{LC}
  \]
  %
  where we'll assume that the source voltage $V_{\text{in}}$ is constant.

  \bigskip

  \begin{enumerate}
  \item[a)] Show that a natural time scale for the oscillatory behaviour of the system is given by $\tau = \sqrt{LC}$ while the damping due to the resitor occurs on a time scale $\tau = \nicefrac{R}{L}$.

  \item[b)] Show that the system admits a single fixed point given by $V_C = V_{\text{in}}$.

  \item[c)] If $V_{\text{in}} = 0$, show that we have
    %
    \[
    \dfrac{d}{dt} \left( \dfrac{1}{2} Li^2 + \dfrac{1}{2}C V_C^2 \right) = - Ri^2
    \]
    %
    where $\nicefrac{Li^2}{2}$ is the magnetic energy in the coil and $\nicefrac{CV_C^2}{2}$ is the electric energy in the condensator.
    From a physical point of view, what does the term $-Ri^2$ represent ?
  \end{enumerate}
\end{problem}

\begin{figure}
  \centering
  \begin{circuitikz}
    \draw (0, 0) to[battery1, invert] (0, 3) to [R, i>^=$i(t)$, l=$R$] (2.5, 3);
    \draw (2.5, 3) to [L, l=$L$] (2.5, 0);
    \draw (2.5, 0) to [C, l=$C$] (0, 0);
  \end{circuitikz}
  \caption{RLC circuit for Problem 2.2}\label{fig: RLC circuit}
\end{figure}

\section{Nonlinear systems}\label{sec: nonlinear systems}

\subsection{List of problems}

\subsubsection*{Conservative systems}

\begin{problem}
  Consider the Duffing oscillator whose equation is given by
  %
  \[
  \ddot{x} + \epsilon x + x^3 = 0
  \]
  %
  where $\epsilon$ can either be positive or negative.

  \bigskip

  \begin{enumerate}
  \item[a)] Show that this system is conservative and give the expression of its Hamiltonian.

  \item[b)] For $\epsilon > 0$, show that the system has a nonlinear center at its origin and sketch the phase portrait.

  \item[c)] For $\epsilon < 0$, show that the system admits three fixed points and classify them.
    Sketch the phase portrait.
  \end{enumerate}
\end{problem}

\bigskip

\begin{problem}
  Consider the \emph{Lotka-Volterra predator-prey model} given by
  %
  \[
  \begin{aligned}
    \dot{R} & = a R - b RF \\
    \dot{F} & = -c F + d RF
  \end{aligned}
  \]
  %
  where $R(t)$ denotes the population of rabbits and $F(t)$ that of foxes.
  We'll assume that all constants are positive as well as $R(t)$ and $F(t)$.

  \begin{enumerate}
  \item[a)] Discuss the biological meaning of all the terms in the model as well as its limitations.

  \item[b)] Introducing a suitable rescaling of the variables $R$ and $F$ as well as time, show that the model can be recast as
    %
    \[
    \begin{aligned}
      \dot{x} & = x - xy \\
      \dot{y} & = - \mu y + \mu xy.
    \end{aligned}
    \]

  \item[c)] Show that the system is conservative.

  \item[d)] Compute the fixed points of the system and classify them.

  \item[e)] Assuming that both $x(0)$ and $y(0)$ are non-zero, show that model predicts cycles in the populations of both species and sketch the phase portrait.

  \item[f)] Sketch the evolution of $x(t)$ and $y(t)$ as a function of time.
    Explain from a biological point of view what you observe.
  \end{enumerate}
\end{problem}

\subsubsection*{Dissipative systems}

\begin{problem}
  Consider once again the Duffing oscillator but let us now include the influence of friction.
  The equation becomes
  %
  \[
  \ddot{x} + \dot{x} - x + x^3 = 0
  \]

  \begin{enumerate}
  \item[a)] Show that the system is no longer conservative.

  \item[b)] Recast the system as two first-order equations, compute the fixed points of the system and classify them.

  \item[c)] In the vicinity of the origin, the stable manifold can be represented as $y = h(x)$.
    Assuming that $h(x) = a x + b x^2 + c x^3$, determine the coefficients of this third-order approximation.
  \end{enumerate}
\end{problem}

\subsubsection*{Modelling a pandemic}

\begin{problem}
  The recent years have been marked by the COVID-19 pandemic.
  Since the begining of the pandemic, numerous models have been proposed to forecast its evolution.
  One of the simplest, dating back to the seminal work of Kermack and McKendrick in the late 1920's and early 1930's, is known as the SIR model.
  It is given by
  %
  \[
  \begin{aligned}
    \dfrac{dS}{dt} & = -\beta \dfrac{SI}{N} \\
    \dfrac{dI}{dt} & = \beta \dfrac{SI}{N} - \gamma I \\
    \dfrac{dR}{dt} & = \gamma I
  \end{aligned}
  \]
  %
  where $N$ is the total population and $S$, $I$ and $R$ denote the number of suceptible (i.e.\ not yet infected), infected and removed (i.e.\ having recovered or died from the disease) persons in the population.
  The parameter $\beta$ characterizes the probability for susceptible individuals to become infected whenever they meet an infectious person.
  The parameter $\gamma$ models how fast an infected individual recovers (or dies) from the disease.

  \begin{enumerate}
  \item[a)] Show that $S(t) + I(t) + R(t)$ is a conserved quantity.

  \item[b)] Introducing a suitable rescaling of the variables and time, show that the model can be recast as
    %
    \[
    \begin{aligned}
      \dot{x} & = -R_0 xy \\
      \dot{y} & = R_0 xy - y \\
      \dot{z} & = y
    \end{aligned}
    \]
    %
    where $R_0$ is known as the \emph{basic reproduction number}.
    From an epidemiological point of view, what does this number represent ?

  \item[c)] The equation for $z$ being decoupled from the rest, the dynamics of $x$ and $y$ indeed forms a two-dimensional system.
    Find and classify all of the fixed points of the system.

  \item[d)] Sketch the nullclines and the phase portrait of the system.

  \item[e)] Let us assume that, as the disease first emerges, most of the population is susceptible (i.e.\ $x_0 \simeq 1$).
    An epidemic is said to occur if $y(t)$ initially increases.
    Under what condition on the basic reproduction number $R_0$ does an epidemic can occur ?

  \item[f)] At the begining of a pandemic, $x(t) \simeq 1$.
    Show that the number of infected individuals initially grows exponentially fast.

  \item[g)] If an outbreak happens, it is considered that \emph{herd immunity} has been achieved once the number of infected individuals starts to decrease.
    What fraction of the population needs to have been infected at one point or another before such a decrease happens ?

  \item[h)] At the begining of the COVID-19, the basic reproduction number was estimated to be approximately equal to 3.
    Given that France has a population of roughly 70 millions inhabitants, how many people would have had to contract the disease before herd immunity had been achieved ?
  \end{enumerate}
\end{problem}

\bigskip

\begin{problem}
  Several strategies have been proposed to mitigate the pandemic.
  Two of the most widely used strategies are social distancing and isolating infected individuals.
  A simple model for the influence of these two strategies is
  %
  \[
  \begin{aligned}
    \dot{x} & = -(R_0 - a) xy \\
    \dot{y} & = (R_0 - a) xy - b y
  \end{aligned}
  \]
  %
  where $a \geq 0$ characterizes by how much we reduced our social interactions and $b \geq 1$ how fast infected individuals are effectively removed from the population.

  \bigskip

  \begin{enumerate}
  \item[a)] Show that social distancing leads to a slower growth of the number of infected individuals at the begining of the pandemic.

  \item[b)] Rescaling time, show that isolation effectively leads to a modified basic reproduction number $R_b < R_0$.
  \end{enumerate}
\end{problem}

\section{Oscillators}

\subsection{The Poincaré-Lindstedt method}

\subsection{Weakly nonlinear self-sustained oscillators}

\subsection{Strongly nonlinear oscillators and relaxation oscillations}

\subsection{List of problems}

\subsubsection*{Poincaré-Lindstedt method}

\begin{problem}
  Consider the conservative Duffing oscillator whose governing equation is given by
  %
  \[
  \ddot{x} + x + \epsilon x^3 = 0 \quad \text{with initial conditions} \quad x(0) = 1, \quad \dot{x}(0) = 0
  \]
  %
  where $\epsilon$ is a small real number.
  The properties of this sytem and its fixed points have already been studied qualitatively in the previous section.
  We now aim for a more quantitative analysis.

  \bigskip

  \begin{enumerate}
  \item[a)] Introducing a strained time variable $\tau = \omega t$ with $\omega = \left( 1 + \epsilon \omega_1 + \cdots\right)$ and assuming that $x(\tau) = x_0(\tau) + \epsilon x_1(\tau) + \cdots$, compute the zeroth and first order approximation of the periodic solution.

  \item[b)] Sketch the evolution of the frequencies in the system as a function of $\epsilon$ (consider both positive and negative values of $\epsilon$).

  \item[c)] Write down a computer program that simulate the system above for different values of $\epsilon$.
    For each value, compare the computer-generated trajectories in phase plane with your Poincaré-Lindstedt approximation.
    Compare the frequencies extracted from your simulation using the Fourier transform with the results from your theoretical analysis.
  \end{enumerate}
  
\end{problem}

\bigskip

\begin{problem}
  When discussing the Poincaré-Lindstedt method for the van der Pol oscillator
  %
  \[
  \ddot{x} + \epsilon \left( x^2 - 1 \right) \dot{x} + x = 0
  \]
  %
  we have seen that, when $\epsilon$ is sufficiently small, the oscillation frequency was given by $\omega = 1 + O(\epsilon^2)$.
  Assuming now that
  %
  \[
  \omega = 1 + \epsilon \omega_1 + \epsilon^2 \omega_2 + O(\epsilon^3)
  \]
  %
  and
  %
  \[
  x(t) = x_0(t) + \epsilon x_1(t) + \epsilon^2 x_2(t) + O(\epsilon^3),
  \]
  %
  show that the frequency is actually given by $\omega = 1 - \dfrac{1}{16} \epsilon^2 + O(\epsilon^3)$.
\end{problem}

\subsubsection*{Weakly nonlinear oscillators}

\begin{problem}[A simple model for the Bénard-von Kàrmàn vortex street]
  Consider the following third-order system
  %
  \[
  \begin{aligned}
    \dot{x} & = \sigma x - y - xz - \alpha yz\\
    \dot{y} & = x + \sigma y - yz + \alpha xz\\
    \dot{z} & = -z + x^2 + y^2.
  \end{aligned}
  \]
  %
  This model arises as a simplified representation of the Bénard-von K\`arm\`an vortex street in the wake of a two-dimensional cylinder at low Reynolds numbers.
  Here, $x$ and $y$ represent the two-degrees of freedom needed to describe the oscillatory behaviour of the flow while $z$ is known as the \emph{shift-mode amplitude}.
  It describes how different the linearly unstable base flow is from the time-averaged mean flow.
  The term $x^2 + y^2$ in the equation for $z$ models the influence of the fluctuation's Reynolds stresses on the mean flow.

  \begin{enumerate}
  \item[a)] Show that this dynamical system is equivariant with respect to the transformation
    %
    \[
    \gamma
    =
    \begin{bmatrix}
      \cos(\theta) & \sin(\theta) & 0 \\
      -\sin(\theta) & \cos(\theta) & 0 \\
      0 & 0 & 1
    \end{bmatrix}
    \]
    %
    where $\gamma$ represents a rotation around the $z$-axis.

  \item[b)] Introducing the complex variable $\eta = x + i y = r e^{i\phi}$, where $r$ is the radius of oscillation and $\phi$ the phase, express the system in the cylindrical coordinates system $(r, \phi, z)$.

  \item[c)] Compute the fixed points in the $(r, z)$ plane.
    You should obtain $(r_0, z_0) = (0, 0)$ and $(r_1, z_1) = (\sqrt{\sigma}, \sigma)$.
    From a physical point of view, to what state does each of these fixed points correspond to ?

  \item[d)] Show that the origin $(r_0, z_0)$ is linearly unstable while the mean flow $(r_1, z_1)$ is linearly stable.

  \item[e)] The equation for the phase reads
    %
    \[
    \dot{\phi} = 1 + \alpha z.
    \]
    %
    How does the oscillation frequency $\omega(t) = 1 + \alpha z(t)$ evolves as the system transitions from its unstable equilibrium towards its limit cycle ?

  \item[f)] Sketch a three-dimensional phase portrait in the coordinate system $(x, y, z)$.
  \end{enumerate}
\end{problem}

\bigskip

\begin{problem}[The Bénard-von Kàrmàn vortex street revisited]
  Let us revisit the simple model for the Bénard-von Kàrmàn vortex street.
  Our goal will now be to derive an amplitude equation for the oscillation and the phase using a weakly nonlinear expansion in the vicinity of the fixed point $(x, y, z) = (0, 0, 0)$.

  \begin{enumerate}
  \item[a)] Introducing the complex variable $\eta = x + i y$, write down the model for $\eta$ and $z$.

  \item[b)] Let us introduce a slow time scale $\tau = \epsilon^2 t$ with $\epsilon^2 = \sigma$.
    Expanding the solution $(\eta, z)$ in the vicinity of $(\eta_0, z_0) = (0, 0)$ as
    %
    \[
    \begin{aligned}
      \eta(t, \tau) & = \epsilon \eta_1(t, \tau) + \epsilon^2 \eta_2(t, \tau) + \epsilon^3 \eta_3(t, \tau) + O(\epsilon^4)\\
      z(t, \tau) & = \epsilon z_1(t, \tau) + \epsilon^2 z_2(t, \tau) + \epsilon^3 z_3(t, \tau) + O(\epsilon^4),
    \end{aligned}
    \]
    %
    show that you obtain the following set of equations for each order
    %
    \[
    \begin{aligned}
      O(\epsilon) : \quad & \dfrac{d\eta_1}{dt} = i \eta_1 \\
      & \dfrac{dz_1}{dt} = - z_1 \\
      O(\epsilon^2) : \quad & \dfrac{d\eta_2}{dt} = i \eta_2 - (1 - i\alpha) \eta_1 z_1 \\
      & \dfrac{dz_2}{dt} = - z_2 + \vert \eta_1 \vert^2 \\
      O(\epsilon^3) : \quad & \dfrac{d\eta_3}{dt} = i \eta_3 + \eta_1 - (1 - i\alpha)(\eta_1 z_2 + \eta_2 z_1) - \dfrac{d\eta_1}{d\tau} \\
      & \dfrac{dz_3}{dt} = - z_3 + \overline{\eta}_1 \eta_2 + \eta_1 \overline{\eta}_2 - \dfrac{dz_1}{d\tau}.
    \end{aligned}
    \]

  \item[c)] The solution for $\eta_1$ is \( \eta_1(t, \tau) = A(\tau) e^{it}\).
    Show that \( z_2(t, \tau) = \vert A(\tau) \vert^2 \left( 1 - e^{-t} \right) \) (assuming $z_2(0, 0) = 0$).

  \item[d)] Introducing these two expressions in the equation for $\eta_3$, show that invoking the Fredholm theorem to kill the resonant terms leads to
    %
    \[
    \dfrac{dA}{d\tau} = A - \left( 1 - i\alpha \right) \vert A \vert^2 A.
    \]
    %
    Expressing the complex amplitude as $A(\tau) = R(\tau) e^{i \phi(\tau)}$, show that we finally have
    %
    \[
    \begin{aligned}
      \dfrac{dR}{d\tau} & = R - R^3 \\
      \dfrac{d\phi}{d\tau} & = \alpha R^2.
    \end{aligned}
    \]

  \item[e)] Given that $\tau = \epsilon^2 t = \sigma t$, show that the solution $\eta(t, \tau)$ to first order is
    %
    \[
    \eta(t, \tau) = \sqrt{\sigma} R(\tau) \exp \left( i\omega(\tau)t \right) + O(\epsilon^2)
    \]
    %
    where $\omega(\tau) = 1 + \alpha \sigma R^2(\tau)$.
    Similarly, assuming that $z_1(0) = 0$, show that the solution for $z(t, \tau)$ up to second order is given by
    %
    \[
    z(t, \tau) = \sigma R^2(\tau) \left( 1 - e^{-t} \right) + O(\epsilon^3).
    \]
    Given that $R(\tau) \to 1$ when $\tau \to +\infty$, are your results regarding the amplitude and frequency of the oscillations and the amplitude of the distortion $z(t)$ consistent with the results from the previous exercise ?
  \end{enumerate}
  
\end{problem}

\subsubsection*{Strongly nonlinear oscillators}

%% \begin{problem}
%%   Consider the FitzHugh-Nagumo model
%%   %
%%   \[
%%   \begin{aligned}
%%     \dot{v} & = v - \dfrac{v^3}{3} - w + I \\
%%     \dot{w} & = \epsilon(v + a - b w).
%%   \end{aligned}
%%   \]
%%   %
%%   This model is a simplified version of the Hodgkin-Huxley model detailing the activation and deactivation dynamics of a spiking neuron.
%%   Here, $v$ denotes the membrane potential while $w$ is known as the recovery variable.
%%   The input $I$ is the current imposed to the membrane.
%%   If $I$ is sufficiently small, no particular long-term dynamics are observed.
%%   However, when $I$ exceeds a critical value, the neuron starts to spike.
%%   These periodic spikes (when $I$ is kept constant above this threshold) can be understood as relaxation oscillations.
%%   Let us try to understand these dynamics using both theoretical and numerical analyses.

%%   \bigskip

%%   \begin{enumerate}
%%   \item[a)] Show that for $a = b = 0$, the FitzHugh-Nagumo model reduces to the classical van der Pol oscillator.
%%   \end{enumerate}
%% \end{problem}

\section{Bifurcations revisited}

\subsection{List of problems}

\begin{problem}[The Bénard-von Kàrmàn problem revisited (again)]

\end{problem}


% --> Resonance and synchronization
\chapter{Resonance and synchronization}


% --> Chaos
\chapter{Chaotic dynamics}
\section{Driven oscillators}

\section{The Lorenz system}

\subsection{List of problems}

\subsubsection*{The Lorenz system}

\begin{problem}[In-depth analysis of the Lorenz system]
  The Lorenz system reads
  %
  \[
  \begin{aligned}
    \dot{x} & = \sigma ( y - x ) \\
    \dot{y} & = x( \rho - z) - y\\
    \dot{z} & = xy - \beta z
  \end{aligned}
  \]
  %
  where $\sigma$, $\rho$ and $\beta$ are all positive parameters.
  In this exercise, you'll dive deep down the rabbit hole and explore this incredible system both from the theoretical and numerical point of view using all of the tools we've discussed so far.

  \paragraph{Simple properties of the Lorenz system}
  %
  \begin{enumerate}
  \item[a)] Show that the Lorenz system is equivariant with respect to the transformation $(x, y, z) \mapsto (-x, -y, z)$.
    What does this properties tell you about the solutions of the system ?

  \item[b)] Show that the $z$-axis is an invariant line, i.e.\ if $x(0) = y(0) = 0$ then $x(t) = y(t) = 0\ \forall t$.

  \item[c)] Show that the Lorenz system is dissipative and that any volume $V$ of initial conditions evolves according to
    %
    \[
    \dot{V} = -\left( \sigma + 1 + b \right) V.
    \]
    %
    Can the Lorenz system exhibit quasi-periodic dynamics ? Explain your reasonning from a geometrical point of view.
    Using a similar argument, explain why the Lorenz system cannot have repelling fixed points (i.e.\ unstable nodes) or repelling orbits.

  \item[d)] Compute the fixed points of the Lorenz system and give their domain of existence.
    In analogy to the Rayleigh-Bénard convection problem, to what type of flow do these fixed points correspond to ?
  \end{enumerate}

  \paragraph{The origin $(x, y, z) = (0, 0, 0)$}
  %
  \begin{enumerate}
    \item[a)] Given the Lyapunov function
    %
    \[
    V(x, y, z, t) = \dfrac{x^2}{\sigma} + y^2 + z^2 \geq 0,
    \]
    %
    show that the origin is the only attractor in the whole phase space (i.e.\ it is globally stable) for $0 \leq \rho \leq 1$.
    (\underline{Hint} : Show that $\dot{V} < 0 \ \forall x, y, z$ and that $V = 0$ for $x = y = z = 0$)

  \item[b)] Compute the Jacobian matrix of the system and show that the origin experiences a bifurcation at $\rho_c = 1$.
    Using a symmetry argument, what type of bifurcation could it be ?

  \item[c)] Let us now start our analysis to determine precisely what type of bifurcation it is.
    For that purpose, let us write $\rho$ as
    %
    \[
    \begin{aligned}
      \rho & = \rho_c + \epsilon \\
      & = 1 + \epsilon.
    \end{aligned}
    \]
    %
    The equation for $y$ then reads
    %
    \[
    \dot{y} = \left( 1 + \epsilon \right) x - y - xz.
    \]
    %
    For $\epsilon = 0$ (i.e.\ $\rho = \rho_c$), the matrix of eigenvectors $\bm{T}$ is given by
    %
    \[
    \bm{T}
    =
    \begin{bmatrix}
      1 & \sigma & 0 \\
      1 & -1 & 0 \\
      0 & 0 & 1
    \end{bmatrix}.
    \]
    %
    Introducing the change of variable $\bm{x} = \bm{T} \bm{u}$, with $\bm{x} = (x, y, z)$ and $\bm{u} = (u, v, w)$, show that the Lorenz system can be recast as
    %
    \[
    \begin{aligned}
      \dot{u} & = \dfrac{\sigma}{1+\sigma} \left( \epsilon - w \right) \left( u + \sigma v \right) \\
      \dot{v} & = -\left( 1 + \sigma \right) v - \dfrac{1}{1 + \sigma} \left( \epsilon r - w \right) \left( u + \sigma v \right) \\
      \dot{w} & = -\beta w + \left( u + \sigma v \right) \left( u - v \right).
    \end{aligned}
    \]

  \item[d)] Let us now supplement the above system with
    %
    \[
    \dot{\epsilon} = 0.
    \]
    %
    The center manifold $W_c$ is given by
    %
    \[
    W_c = \left\{ (u, v, w, \epsilon) : v = h_1(u, \epsilon), w = h_2(u, \epsilon), h_i(0, 0) = 0, Dh_i(0, 0) = 0 \right\}
    \]
    %
    where $Dh_i$ denotes the Jacobian matrix of the nonlinear parametrization $h_i(u, \epsilon)$.
    Assume now that
    %
    \[
    \begin{aligned}
      h_1(u, \epsilon) = a_{00} + a_{10} u + a_{01} \epsilon + a_{20} u^2 + a_{11} u\epsilon + a_{22} \epsilon^2 + \cdots \\
      h_2(u, \epsilon) = b_{00} + b_{10} u + b_{01} \epsilon + b_{20} u^2 + b_{11} u\epsilon + b_{22} \epsilon^2 + \cdots
    \end{aligned}
    \]
    %
    show that the condition $h_i(0, 0)$ implies
    %
    \[
    a_{00} = b_{00} = 0.
    \]
    %
    Similarly, show that $Dh_i(0, 0) = 0$ implies
    %
    \[
    \begin{aligned}
      a_{10} & = a_{01} = 0 \\
      b_{10} & = b_{01} = 0.
    \end{aligned}
    \]

  \item[e)] Our parametrizations reduce to
    %
    \[
    \begin{aligned}
      h_1(u, \epsilon) & = a_{20} u^2 + a_{11} u \epsilon + a_{22} \epsilon^2 \\
      h_2(u, \epsilon) & = b_{20} u^2 + b_{11} u \epsilon + b_{22} \epsilon^2.
    \end{aligned}
    \]
    %
    Using the fact that
    %
    \[
    \dfrac{dh_i}{dt} = \dfrac{\partial h_i}{\partial u} \dot{u} + \dfrac{\partial h_i}{\partial \epsilon} \dot{\epsilon},
    \]
    %
    show that the governing equation for $u$ can be expressed as
    %
    \[
    \dot{u} = \mu u + \lambda u^3
    \]
    %
    where $\mu$ and $\lambda$ are to be determined.

  \item[f)] Given the expressions of $\mu$ and $\lambda$ obtained at the previous question, conclude about the nature of the bifurcation happening at $\rho = 1$.
  \end{enumerate}


  \paragraph{The fixed points $C^+$ and $C^-$}

  Let us now turn our attention to the two symmetric fixed points newly created hereafter denoted as
  %
  \[
  \begin{aligned}
    C^+ & = (\sqrt{\beta(\rho -1)}, \sqrt{(\beta(\rho-1)}, \rho-1) \\
    C^- & = (-\sqrt{\beta(\rho -1)}, -\sqrt{(\beta(\rho-1)}, \rho-1).
  \end{aligned}
  \]

  \begin{enumerate}
  \item[a)] Show that the characteristic polynomial of the Jacobian matrix at $C^+$, $C^-$ is
    %
    \[
    \lambda^3 + \left( \sigma + \beta + 1 \right) \lambda^2 + \left( \rho + \sigma \right) \beta \lambda + 2 \beta \sigma( \rho - 1) = 0.
    \]
    %
    (\underline{Hint} : Use the fact that, for a $3 \times 3$ matrix $\bm{A}$, its characteristic polynomial is given by $\lambda^3 - \text{tr}(\bm{A})\lambda^2 - \dfrac{1}{2}\left( \text{tr}(\bm{A})^2 - \text{tr}(\bm{A}^2) \right) \lambda - \det(\bm{A}) = 0$.)

  \item[b)] As $\rho$ is increased, $C^+$ and $C^-$ eventually experience a subcritical Hopf bifurcation.
    Assuming that $\lambda = \pm i \omega$, show that this bifurcation happens for
    %
    \[
    \rho_H = \sigma \dfrac{\sigma + \beta + 3}{\sigma - \beta - 1}.
    \]
    %
    In doing so, why do we need to assume $\sigma > \beta + 1$ ?

  \item[c)] The bifurcation being subcritical, there exists no stable limit cycle in the direct vicinity of $C^+$ and $C^-$.
    As discussed in \textsection~\ref{??}, the trajectories nonetheless remain bounded in phase space.
    Show that there exist a ellipsoidal region $\Omega$ of the form
    %
    \[
    \rho x^2 + \sigma y^2 + \sigma(z - 2\rho)^2 \leq C
    \]
    %
    where $C$ is finite positive number such that all trajectories of the Lorenz system eventually enter $\Omega$ and stay in there forever.

  \item[d)] Given all of the results obtained so far, sketch the current status of the bifurcation diagram of the Lorenz system.
  \end{enumerate}

  \paragraph{Chaotic dynamics}

  Let us now explore the chaotic dynamics of the Lorenz system numerically.
  For that purpose, we will hereafter consider the classical parameters
  %
  \[
  \rho = 28, \quad \sigma = 1, \quad \text{and} \quad \beta = \dfrac{8}{3}.
  \]

  \begin{enumerate}
  \item[a)] Using your favorite programming language, simulate the Lorenz system for this set of parameters and plot its phase portrait.
    You can discard any initial transients.

  \item[b)] Having computed a sufficiently long time-series of $z(t)$, extract the successive maxima $z_0, z_1, z_2, \cdots$ and plot the Lorenz map $z_{k+1} = f(z_k)$ along with the bisectrice $z_{k+1} = z_k$.
  \end{enumerate}
\end{problem}

\subsubsection*{Other Lorenz-like systems}

\begin{problem}[Laser model]
  As discussed in \cite{??} and \cite{??}, the Maxwell-Block equations for a laser are
  %
  \[
  \begin{aligned}
    \dot{E} & = \kappa \left( P - E \right) \\
    \dot{P} & = \gamma_1 \left( ED - p \right) \\
    \dot{D} & = \gamma_2 \left( \lambda + 1 - D - \lambda EP \right)
  \end{aligned}
  \]
  %
  where \ldots
  
  \begin{enumerate}
  \item[a)] Find a change of variables that transform this system of equations into the Lorenz system.
  \item[b)] Show that the supercritical pitchfork happening for $\rho = 1$ in the Lorenz system corresponds to the non-lasing state ($E^* = 0$) in the Maxwell-Block loosing its stability.
  \end{enumerate}
\end{problem}

\bigskip

\begin{problem}[The Malkus waterwheel]
  As discussed in \cite{??}, the equations governing the dynamics of the Malkus waterwheel shown in figure~\ref{fig: malkus waterwheel} are given by
  %
  \[
  \begin{aligned}
    \dot{a} & = \omega b - K a \\
    \dot{b} & = -\omega a + q - Kb \\
    \dot{\omega} & = -\dfrac{\nu}{I} \omega + \dfrac{\pi g r}{I} a,
  \end{aligned}
  \]
  %
  where \ldots .

  \begin{enumerate}
  \item[a)] Find a change of variables transforming the waterwheel equations into the Lorenz system.
  \item[b)] Show that, for the water wheel, $\beta = 1$ in the corresponding Lorenz system.
  \item[c)] As shown when discussing the Lorenz system, a supercritical pitchfork bifurcation occurs at $\rho = 1$.
    To what kind of motion do the newly created fixed points correspond to for the waterwheel ?
  \end{enumerate}
\end{problem}

\section{Unimodal maps}

\section{A primer in fractal geometry}

\subsection{List of problems}

\subsubsection*{Self-similar fractals}

\def\trianglewidth{.22\textwidth}%
\pgfdeclarelindenmayersystem{Sierpinski triangle}{
    \symbol{X}{\pgflsystemdrawforward}
    \symbol{Y}{\pgflsystemdrawforward}
    \rule{X -> X-Y+X+Y-X}
    \rule{Y -> YY}
}%

\def\kochcurve{.333\textwidth}%
\pgfdeclarelindenmayersystem{Koch curve}{
  \symbol{X}{\pgflsystemdrawforward}
  \rule{X -> X + X -- X + X}
}

\pgfdeclarelindenmayersystem{three half curve}{
  \symbol{X}{\pgflsystemdrawforward}
  \rule{X -> X + X - X - XX + X + X - X}
}

\begin{figure}
  \centering
  \subfigure[$n=0$]{
    \begin{tikzpicture}
      \draw[thick] (0,0) l-system
           [l-system={Koch curve, axiom=X, step=\kochcurve, order=0, angle=60}];
    \end{tikzpicture}    
  }
  \subfigure[$n=1$]{
    \begin{tikzpicture}
      \draw[thick] (0,0) l-system
           [l-system={Koch curve, axiom=X, step=\kochcurve/3, order=1, angle=60}];
    \end{tikzpicture}    
  }\\

  \subfigure[$n=2$]{
    \begin{tikzpicture}
      \draw[thick] (0,0) l-system
           [l-system={Koch curve, axiom=X, step=\kochcurve/9, order=2, angle=60}];
    \end{tikzpicture}    
  }
  \subfigure[$n=3$]{
    \begin{tikzpicture}
      \draw[thick] (0,0) l-system
           [l-system={Koch curve, axiom=X, step=\kochcurve/27, order=3, angle=60}];
    \end{tikzpicture}    
  }

  \caption{First four iterations for the construction of the Koch curve.}\label{fig: koch curve}
\end{figure}

\begin{problem}
  Consider the fractal object known as the Koch curve.
  The first four iterations to construct this curve are shown in figure~\ref{fig: koch curve}.
  It can be generated by the same four affine transformations again and again.

  \begin{enumerate}
  \item[a)] Find all four affine transformations needed to construct this object and show that they are contractive (i.e.\ all of their eigenvalues are inside the unit circle).

  \item[b)] Show that its perimeter goes to infinity as the number of iterations $n \to \infty$.
  \end{enumerate}
\end{problem}

\bigskip

\begin{figure}
  \centering
  \subfigure[$n=0$]{
    \begin{tikzpicture}
      \draw[thick] (0,0) l-system
           [l-system={three half curve, axiom=X, step=\kochcurve, order=0, angle=90}];
    \end{tikzpicture}    
  }
  \subfigure[$n=1$]{
    \begin{tikzpicture}
      \draw[thick] (0,0) l-system
           [l-system={three half curve, axiom=X, step=\kochcurve/4, order=1, angle=90}];
    \end{tikzpicture}    
  }\\

  \subfigure[$n=2$]{
    \begin{tikzpicture}
      \draw[thick] (0,0) l-system
           [l-system={three half curve, axiom=X, step=\kochcurve/16, order=2, angle=90}];
    \end{tikzpicture}    
  }
  \subfigure[$n=3$]{
    \begin{tikzpicture}
      \draw[thick] (0,0) l-system
           [l-system={three half curve, axiom=X, step=\kochcurve/64, order=3, angle=90}];
    \end{tikzpicture}    
  }

  \caption{First four iterations for the construction of the $\dfrac{3}{2}$ curve.}\label{fig: 3/2 curve}
\end{figure}

\begin{problem}
  Consider the curve shown in figute~\ref{fig: 3/2 curve}.
  This curve is simply known as the $\dfrac{3}{2}$ curve.
  Why is it call that way ?
\end{problem}

\bigskip

\begin{figure}
  \centering
  \subfigure[$n=0$]{
    \begin{tikzpicture}
      \fill [black] (0,0) -- ++(0:\trianglewidth) -- ++(120:\trianglewidth) -- cycle;
      \draw [draw=none] (0,0) l-system
            [l-system={Sierpinski triangle, axiom=X, step=\trianglewidth, order=0, angle=-120},fill=white];
    \end{tikzpicture}    
  }%
  \hfill
  \subfigure[$n=1$]{
    \begin{tikzpicture}
      \fill [black] (0,0) -- ++(0:\trianglewidth) -- ++(120:\trianglewidth) -- cycle;
      \draw [draw=none] (0,0) l-system
            [l-system={Sierpinski triangle, axiom=X, step=\trianglewidth/2, order=1, angle=-120},fill=white];
    \end{tikzpicture}    
  }%
  \hfill
  \subfigure[$n=2$]{
    \begin{tikzpicture}
      \fill [black] (0,0) -- ++(0:\trianglewidth) -- ++(120:\trianglewidth) -- cycle;
      \draw [draw=none] (0,0) l-system
            [l-system={Sierpinski triangle, axiom=X, step=\trianglewidth/4, order=2, angle=-120},fill=white];
    \end{tikzpicture}    
  }%
  \hfill
  \subfigure[$n=3$]{
    \begin{tikzpicture}
      \fill [black] (0,0) -- ++(0:\trianglewidth) -- ++(120:\trianglewidth) -- cycle;
      \draw [draw=none] (0,0) l-system
            [l-system={Sierpinski triangle, axiom=X, step=\trianglewidth/8, order=3, angle=-120},fill=white];
    \end{tikzpicture}    
  }

  \caption{First four iterations for the construction of the Sierpinsky triangle.}\label{fig: sierpinsky triangle}
\end{figure}

\begin{problem}
  Consider the fractal object known as the Sierpinsky triangle whose first four iterations are shown in figure~\ref{fig: sierpinsky triangle}.
  This fractal can be generated by applying repeatdly the same three affine transformations.

  \begin{enumerate}
  \item[a)] Find all three affine transformations needed to construct this figure and show that they are are contractive (i.e.\ all of their eigenvalues are inside the unit circle).

  \item[b)] Using simple geometric arguments, show that the perimeter of the Sierpinsky triangle goes to infinity as $n \to \infty$ (where $n$ is the number of iterations) while its area goes to zero.

  \item[c)] What is the Haussdorf dimension of this fractal ?
  \end{enumerate}
\end{problem}

\subsubsection*{The Mandelbrot set}

\begin{problem}[Logistic map and the Mandelbrot set]
  As discussed in \textsection~\ref{}, the logistic map
  %
  \[
  x_{k+1} = \mu x_k \left( 1 - x_k \right), \quad \text{with} \quad \mu \in \left[0, 4 \right]
  \]
  %
  is actually the Mandelbrot set generating system $z_{k+1} = z_k^2 + c$ in disguise.
  Show that it is indeed the case. (\underline{Hint} : use a simple linear transformation $z = ax + b$)
\end{problem}

\bigskip

\begin{problem}[]
  Prove by induction that any point $c \in \mathbb{C}$ such that $\vert c \vert > 2$ is not in the Mandelbrot set (i.e.\ the iteration $z_{k+1} = z_k^2 + c$ eventually diverges).
\end{problem}

\bigskip

\begin{problem}[Escape radius]
  Prove by induction that if $\vert z_k \vert > 2$ for any point $c \in \mathbb{C}$ such that $\vert c \vert < 2$, then $c$ is not in the Mandelbrot set.
\end{problem}

\bigskip

\begin{problem}[Approximating $\pi$ in a terribly inefficient way]
  $\pi$ has the tendancy of appearing in surprising ways in numerous problems where, at first, it seems like it shouldn't.
  Let us consider $c = -\dfrac{3}{4} + i \epsilon$ where $\epsilon$ is a small positive real number.
  Note that $c = -\dfrac{3}{4}$ is in the Mandelbrot set.
  It is thus expected that, as $\epsilon \to 0$, an increasing number of steps is required for the iteration to reach the escape radius $\vert z_k \vert > 2$.
  Let us denote this number by $N(\epsilon)$.
  Write down a small program that compute $N(\epsilon)$ for $\epsilon = 1, 0.1, 0.01, 0.001, \cdots$ and see what happens !

\end{problem}

\section{Back to strange attractors}


% --> Other types of dynamical systems.
\chapter{Other types of dynamical systems}

\section{Hyper-chaotic systems}

\section{Cellular automata}

\section{Stochastic dynamical systems}

\section{Spatio-temporal systems}


%----------------------------------------------------------------------------------------
%	PART TWO : MACHINE LEARNING
%----------------------------------------------------------------------------------------

\part{Introduction to Machine Learning}

%----------------------------------------------------------------------------------------
%	PART THREE : DATA-DRIVEN METHODS FOR DYNAMICAL SYSTEMS
%----------------------------------------------------------------------------------------

\part{Data-driven methods for dynamical systems}



%----------------------------------------------------------------------------------------
%	CHAPTER 1
%----------------------------------------------------------------------------------------

%% \chapterimage{chapter_head_2.pdf} % Chapter heading image

%% \chapter{Text Chapter}

%% \section{Paragraphs of Text}\index{Paragraphs of Text}

%% \lipsum[1-7] % Dummy text

%% %------------------------------------------------

%% \section{Citation}\index{Citation}

%% This statement requires citation \cite{article_key}; this one is more specific \cite[162]{book_key}.

%% %------------------------------------------------

%% \section{Lists}\index{Lists}

%% Lists are useful to present information in a concise and/or ordered way\footnote{Footnote example...}.

%% \subsection{Numbered List}\index{Lists!Numbered List}

%% \begin{enumerate}
%% \item The first item
%% \item The second item
%% \item The third item
%% \end{enumerate}

%% \subsection{Bullet Points}\index{Lists!Bullet Points}

%% \begin{itemize}
%% \item The first item
%% \item The second item
%% \item The third item
%% \end{itemize}

%% \subsection{Descriptions and Definitions}\index{Lists!Descriptions and Definitions}

%% \begin{description}
%% \item[Name] Description
%% \item[Word] Definition
%% \item[Comment] Elaboration
%% \end{description}

%% %----------------------------------------------------------------------------------------
%% %	CHAPTER 2
%% %----------------------------------------------------------------------------------------

%% \chapter{In-text Elements}

%% \section{Theorems}\index{Theorems}

%% This is an example of theorems.

%% \subsection{Several equations}\index{Theorems!Several Equations}
%% This is a theorem consisting of several equations.

%% \begin{theorem}[Name of the theorem]
%% In $E=\mathbb{R}^n$ all norms are equivalent. It has the properties:
%% \begin{align}
%% & \big| ||\mathbf{x}|| - ||\mathbf{y}|| \big|\leq || \mathbf{x}- \mathbf{y}||\\
%% &  ||\sum_{i=1}^n\mathbf{x}_i||\leq \sum_{i=1}^n||\mathbf{x}_i||\quad\text{where $n$ is a finite integer}
%% \end{align}
%% \end{theorem}

%% \subsection{Single Line}\index{Theorems!Single Line}
%% This is a theorem consisting of just one line.

%% \begin{theorem}
%% A set $\mathcal{D}(G)$ in dense in $L^2(G)$, $|\cdot|_0$. 
%% \end{theorem}

%% %------------------------------------------------

%% \section{Definitions}\index{Definitions}

%% This is an example of a definition. A definition could be mathematical or it could define a concept.

%% \begin{definition}[Definition name]
%% Given a vector space $E$, a norm on $E$ is an application, denoted $||\cdot||$, $E$ in $\mathbb{R}^+=[0,+\infty[$ such that:
%% \begin{align}
%% & ||\mathbf{x}||=0\ \Rightarrow\ \mathbf{x}=\mathbf{0}\\
%% & ||\lambda \mathbf{x}||=|\lambda|\cdot ||\mathbf{x}||\\
%% & ||\mathbf{x}+\mathbf{y}||\leq ||\mathbf{x}||+||\mathbf{y}||
%% \end{align}
%% \end{definition}

%% %------------------------------------------------

%% \section{Notations}\index{Notations}

%% \begin{notation}
%% Given an open subset $G$ of $\mathbb{R}^n$, the set of functions $\varphi$ are:
%% \begin{enumerate}
%% \item Bounded support $G$;
%% \item Infinitely differentiable;
%% \end{enumerate}
%% a vector space is denoted by $\mathcal{D}(G)$. 
%% \end{notation}

%% %------------------------------------------------

%% \section{Remarks}\index{Remarks}

%% This is an example of a remark.

%% \begin{remark}
%% The concepts presented here are now in conventional employment in mathematics. Vector spaces are taken over the field $\mathbb{K}=\mathbb{R}$, however, established properties are easily extended to $\mathbb{K}=\mathbb{C}$.
%% \end{remark}

%% %------------------------------------------------

%% \section{Corollaries}\index{Corollaries}

%% This is an example of a corollary.

%% \begin{corollary}[Corollary name]
%% The concepts presented here are now in conventional employment in mathematics. Vector spaces are taken over the field $\mathbb{K}=\mathbb{R}$, however, established properties are easily extended to $\mathbb{K}=\mathbb{C}$.
%% \end{corollary}

%% %------------------------------------------------

%% \section{Propositions}\index{Propositions}

%% This is an example of propositions.

%% \subsection{Several equations}\index{Propositions!Several Equations}

%% \begin{proposition}[Proposition name]
%% It has the properties:
%% \begin{align}
%% & \big| ||\mathbf{x}|| - ||\mathbf{y}|| \big|\leq || \mathbf{x}- \mathbf{y}||\\
%% &  ||\sum_{i=1}^n\mathbf{x}_i||\leq \sum_{i=1}^n||\mathbf{x}_i||\quad\text{where $n$ is a finite integer}
%% \end{align}
%% \end{proposition}

%% \subsection{Single Line}\index{Propositions!Single Line}

%% \begin{proposition} 
%% Let $f,g\in L^2(G)$; if $\forall \varphi\in\mathcal{D}(G)$, $(f,\varphi)_0=(g,\varphi)_0$ then $f = g$. 
%% \end{proposition}

%% %------------------------------------------------

%% \section{Examples}\index{Examples}

%% This is an example of examples.

%% \subsection{Equation and Text}\index{Examples!Equation and Text}

%% \begin{example}
%% Let $G=\{x\in\mathbb{R}^2:|x|<3\}$ and denoted by: $x^0=(1,1)$; consider the function:
%% \begin{equation}
%% f(x)=\left\{\begin{aligned} & \mathrm{e}^{|x|} & & \text{si $|x-x^0|\leq 1/2$}\\
%% & 0 & & \text{si $|x-x^0|> 1/2$}\end{aligned}\right.
%% \end{equation}
%% The function $f$ has bounded support, we can take $A=\{x\in\mathbb{R}^2:|x-x^0|\leq 1/2+\epsilon\}$ for all $\epsilon\in\intoo{0}{5/2-\sqrt{2}}$.
%% \end{example}

%% \subsection{Paragraph of Text}\index{Examples!Paragraph of Text}

%% \begin{example}[Example name]
%% \lipsum[2]
%% \end{example}

%% %------------------------------------------------

%% \section{Exercises}\index{Exercises}

%% This is an example of an exercise.

%% \begin{exercise}
%% This is a good place to ask a question to test learning progress or further cement ideas into students' minds.
%% \end{exercise}

%% %------------------------------------------------

%% \section{Problems}\index{Problems}

%% \begin{problem}
%% What is the average airspeed velocity of an unladen swallow?
%% \end{problem}

%% %------------------------------------------------

%% \section{Vocabulary}\index{Vocabulary}

%% Define a word to improve a students' vocabulary.

%% \begin{vocabulary}[Word]
%% Definition of word.
%% \end{vocabulary}

%% %----------------------------------------------------------------------------------------
%% %	PART
%% %----------------------------------------------------------------------------------------

%% \part{Part Two}

%% %----------------------------------------------------------------------------------------
%% %	CHAPTER 3
%% %----------------------------------------------------------------------------------------

%% \chapterimage{chapter_head_1.pdf} % Chapter heading image

%% \chapter{Presenting Information}

%% \section{Table}\index{Table}

%% \begin{table}[h]
%% \centering
%% \begin{tabular}{l l l}
%% \toprule
%% \textbf{Treatments} & \textbf{Response 1} & \textbf{Response 2}\\
%% \midrule
%% Treatment 1 & 0.0003262 & 0.562 \\
%% Treatment 2 & 0.0015681 & 0.910 \\
%% Treatment 3 & 0.0009271 & 0.296 \\
%% \bottomrule
%% \end{tabular}
%% \caption{Table caption}
%% \label{tab:example} % Unique label used for referencing the table in-text
%% %\addcontentsline{toc}{table}{Table \ref{tab:example}} % Uncomment to add the table to the table of contents
%% \end{table}

%% Referencing Table \ref{tab:example} in-text automatically.

%% %------------------------------------------------

%% \section{Figure}\index{Figure}

%% \begin{figure}[h]
%% \centering\includegraphics[scale=0.5]{placeholder.jpg}
%% \caption{Figure caption}
%% \label{fig:placeholder} % Unique label used for referencing the figure in-text
%% %\addcontentsline{toc}{figure}{Figure \ref{fig:placeholder}} % Uncomment to add the figure to the table of contents
%% \end{figure}

%% Referencing Figure \ref{fig:placeholder} in-text automatically.

%----------------------------------------------------------------------------------------
%	BIBLIOGRAPHY
%----------------------------------------------------------------------------------------

\chapter*{Bibliography}
\addcontentsline{toc}{chapter}{\textcolor{gray}{Bibliography}} % Add a Bibliography heading to the table of contents

%------------------------------------------------

\section*{Articles}
\addcontentsline{toc}{section}{Articles}
\printbibliography[heading=bibempty,type=article]

%------------------------------------------------

\section*{Books}
\addcontentsline{toc}{section}{Books}
\printbibliography[heading=bibempty,type=book]

%----------------------------------------------------------------------------------------
%	INDEX
%----------------------------------------------------------------------------------------

\cleardoublepage % Make sure the index starts on an odd (right side) page
\phantomsection
\setlength{\columnsep}{0.75cm} % Space between the 2 columns of the index
\addcontentsline{toc}{chapter}{\textcolor{gray}{Index}} % Add an Index heading to the table of contents
\printindex % Output the index

%----------------------------------------------------------------------------------------

\end{document}
